% introduction.tex
Generating novel protein sequences that fold correctly and perform a desired function is a
central goal of computational protein engineering. Several classes of methods now address this
problem. Deep generative models, including variational autoencoders~\cite{hawkinsHookerVAE2021},
autoregressive language models~\cite{madaniProGen2023}, and diffusion
models~\cite{alamdariProteinGeneration2023}, learn sequence distributions from large databases
and can produce diverse candidates, but require substantial training data and GPU resources.
Structure-conditioned inverse folding methods~\cite{dauparasProteinMPNN2022} design sequences for
a given backbone but require a target structure as input. Protein language
models~\cite{esmfold2022} trained on evolutionary databases can score or sample sequences, but
conditioning them on a specific functional property typically requires fine-tuning on labeled
data. Profile hidden Markov models (HMMs)~\cite{hmmer3} provide a standard training-free model
of family alignments. Sequence weights can shift their position-specific emission marginals
toward a designated subset, making a weighted profile HMM a direct comparator for the present
method. Profile HMMs do not model correlations between match-state emissions, whereas
stochastic attention samples a continuous representation of complete stored sequences. We
therefore compared the two methods under the same alignments, designation labels, and multiplicity ratios. For the
scarce-data regime considered here (families with tens to hundreds of members, typical of Pfam
seed alignments), models trained on the target family alone have limited training signal.
Pretrained models are not restricted in this way, since they draw on sequence data beyond the
target family.

Stochastic attention (SA) was developed for the scarce-data regime where traditional methods struggle due to limited training data~\cite{varnerSAProtein2026}. Stochastic attention treats the energy function of a modern Hopfield network, constructed directly from a small family alignment, as a Boltzmann density and samples from it using Langevin dynamics. The method requires no training, runs on a laptop, and produces sequences that preserve family-level amino acid composition, per-position conservation, and predicted three-dimensional fold. A key
limitation, however, is that SA treats all sequences in the alignment equally. Every family
member contributes the same weight to the energy landscape, so the sampler explores the full
family distribution without preference for any functional subset. In practice, experimentalists
often want something more specific: not just a plausible Kunitz domain, but a Kunitz domain
that \emph{inhibits trypsin}; not just a plausible SH3 domain, but one that binds a
particular proline-rich motif. The members of a family alignment share a fold but differ in
function, and the residues that carry the difference occupy only a few positions. A sampler
that reproduces the family distribution faithfully therefore reproduces its variation at those
positions, producing the functionally relevant residue only as often as it appears in the
alignment. This raises the question of how to condition the generative process on a functional
subset without retraining.

In this study, we answered that question with multiplicity weighting. Stochastic attention
encodes each aligned sequence as a column in a memory matrix. At each step the sampler
computes attention weights over these stored sequences and moves toward their weighted average.
A bias added to the attention scores makes it spend more time near selected sequences. We
assigned a multiplicity weight $r_k > 0$ to each stored sequence and added $\log r_k$ to its
attention score before the softmax normalization. All designated sequences share one weight and all
background sequences share another, so a single number, the multiplicity ratio $\rho$ between
the two weights, controls the strength of the bias: $\rho = 1$ recovers standard SA (no
preference), while increasing $\rho$ shifts generation toward the designated subset. 
The method is agnostic to what the designated subset represents (binding, thermostability, organism of origin); it simply biases the sampler toward the sequences the user points to. 
The same multiplicity-weighted mechanism has also been applied to synthetic patient generation
from small longitudinal clinical cohorts~\cite{varnerSAPatient2026}, so the conditioning is not
specific to sequence data. We compared multiplicity weighting against the obvious alternative,
hard curation, which builds the memory matrix from the designated sequences alone and leaves
the rest of the family out. Those discarded sequences are what show which
positions the fold holds fixed and which are free to vary, so dropping them leaves the sampler
with far fewer examples to work from. We tested both on five
Pfam protein families (Kunitz, SH3, WW, Homeobox, and Forkhead domains) spanning a range of
family sizes ($K = 55$--$420$) and designated-subset geometries. Hard curation retained the
selected marker completely from as few as three designated inputs. Those inputs were chosen
because they carry the marker, so recovering it demonstrated control over the sampler and not
that the generated sequences had the function the marker stands for. The multiplicity weights,
by contrast, enter the sampler as a simple logit bias, preserving the analytic score function and
$\mathcal{O}(dK)$ per-step cost of the original method. The conditioning has two stages, and only the
first was exact. Inside the sampler, the mean attention mass on the designated patterns matched
the fraction that $\rho$ sets to within \AttnMaxDevPts{} percentage points at worst
(\AttnMaxDevFamily{}, $\rho=\AttnMaxDevRho$) across the canonical replicated sweeps. Turning
those continuous states back into discrete sequences cost part of that conditioning: after
finite-temperature sampling, reconstruction from unit-normalized principal component analysis
(PCA) coordinates, and argmax decoding, the fraction of generated sequences carrying the marker
could fall below the attention mass the sampler had placed on the designated set. We termed that
shortfall the calibration gap and examined it against a simple geometric quantity, the
separation index $S$, which measures how well PCA separates designated from background
sequences. Multiplicity weighting therefore offers a practical route to
conditioning generative models in the scarce-data regime, and the calibration gap is the
quantity anyone applying such conditioning should measure before trusting it.
